\section{\texorpdfstring{$\mathbb P^3$}{P3} の量子化と関数体}
\label{sec:p3-quantum-models}\label{sec:p3-models}

Pym \cite[定理1.1、表1、第3節]{Pym2015} は、四変数多項式環の
次数付き Calabi--Yau 変形の六成分と一般元の関係式を記述した。
対応する射影 Poisson 成分は
$L(1,1,1,1)$、$L(1,1,2)$、$R(2,2)$、$R(1,3)$、$S(2,3)$、$E(3)$ である。
それぞれ skew-polynomial 環、二次の Ore 拡大、四次元 Sklyanin 代数、
三次元 Sklyanin 代数の中心拡大、可換多項式環の微分 Ore 拡大、
および二次元非可換 Lie 代数に対する CGG 積によって量子化される。
以下では $Q_0(A)$ を次数付き分数環の次数零部分とする。

\subsection{引き戻し成分の微分 Ore 表示}

$F_i\in\mathbb C[x_1,x_2,x_3]_2$ とし、
$\delta=\sum_{i=1}^3F_i\partial_{x_i}$ と置く。
Pym の $S(2,3)$ 型は $\operatorname{div}\delta=0$ の場合に対応し、
量子化は
\[
 A_h=\mathbb C[h,x_1,x_2,x_3][x_0;h\delta]
\]
で与えられる。$x_0,x_1,x_2,x_3$ の次数を1とする。
可換係数環の非零元を Ore 局所化し、
\[
 v=x_1/x_3,\qquad w=x_2/x_3,\qquad T=x_3^{-1}x_0,
 \qquad f_i(v,w)=F_i(v,w,1)
\]
と置けば
\[
 Q_0(A_h)=\operatorname{Frac}\bigl(\mathbb C(h)(v,w)[T;hD]\bigr),
 \qquad
 D=(f_1-vf_3)\partial_v+(f_2-wf_3)\partial_w.
\]
実際、$[T,v]=h x_3^{-1}\delta(v)=h(f_1-vf_3)$ であり、$w$ も同様である。
局所化した次数付き環は、この次数零の環に $x_3^{\pm1}$ を加えて得られる。
共役は $x_3Tx_3^{-1}=T-hf_3$、$x_3v=vx_3$、$x_3w=wx_3$ なので、
その次数零の分数環は表示した微分 Ore 環の分数環に一致する。

\subsection{例外成分の Weyl 関数体}
\label{subsec:exceptional-weyl}

Pym \cite[第3.6節、補題3.3]{Pym2015} の二つの導分を
\[
 X=-\frac54x_0\partial_{x_0}-\frac14x_1\partial_{x_1}
       +\frac34x_2\partial_{x_2}+\frac74x_3\partial_{x_3},
 \qquad
 Y=4(x_0\partial_{x_1}+x_1\partial_{x_2}+x_2\partial_{x_3})
\]
とする。$[Y,X]=Y$ であり、積を
\begin{equation}\label{eq:exceptional-cgg}
 f\star g=\sum_{k\ge0}h^kY^k(f)\binom{X}{k}(g),
 \qquad \binom Xk=\frac{X(X-1)\cdots(X-k+1)}{k!}
\end{equation}
で定める。この積の一次の交換子を半古典括弧の符号規約とする。
$Y$ は局所冪零だから、この式は各 $f,g$ に対して有限和となる。

\begin{proposition}\label{prop:exceptional-weyl}
式\eqref{eq:exceptional-cgg}による例外型量子化の次数零の関数体は、
$h\ne0$ の各ファイバーにおいて
\[
 \operatorname{Frac}\bigl(\mathbb C(c)\langle r,P\rangle/(Pr-rP-h)\bigr)
\]
と同型である。従ってその中心は $\mathbb C(c)$ である。
\end{proposition}
\begin{proof}
$Y(x_0)=0$ であり、$x_0$ は $X$ の固有ベクトルだから、$x_0$ は正規正則元である。
具体的に $a=-5/4$ と置くと
$f\star x_0=x_0\star(1+hY)^a f$ であり、$(1+hY)^a$ は局所冪零性により
可逆な線形写像である。従って $x_0$ を Ore 局所化できる。
同じ式は可換環 $\mathbb C[x_0^{\pm1},x_1,x_2,x_3]$ 上で有限和のまま定義され、
その次数零部分の可換座標 $z_i=x_i/x_0$ 上の導分は
\[
 \overline X=z_1\partial_{z_1}+2z_2\partial_{z_2}+3z_3\partial_{z_3},
 \qquad
 \overline Y=4(\partial_{z_1}+z_1\partial_{z_2}+z_2\partial_{z_3})
\]
となる。ここで $z_i$ は積を変形する前のベクトル空間の座標である。
\[
 s=\frac{z_1}{4},\qquad
 p=z_2-\frac{z_1^2}{2},\qquad
 q=z_3-z_1z_2+\frac{z_1^3}{3}
\]
と置くと、可換環として $\mathbb C[z_1,z_2,z_3]=\mathbb C[s,p,q]$ であり、
\[
 \overline Y=\partial_s,\qquad
 \overline X=s\partial_s+2p\partial_p+3q\partial_q.
\]
従って式\eqref{eq:exceptional-cgg}から厳密に
\[
 [s,p]_\star=2hp,\qquad [s,q]_\star=3hq,\qquad [p,q]_\star=0
\]
を得る。$\star$ 冪 $s^{\star k}$ は通常の $s^k$ と、$s$ に関する低次項の和である。
従って $p^iq^js^{\star k}$ は基底をなし、チャートの代数は
\[
 \mathbb C[h,p,q][s;h(2p\partial_p+3q\partial_q)]
\]
に同型である。

分数環で
\[
 r=q/p,\qquad c=q^2/p^3,\qquad P=r^{-1}s
\]
と置く。$p,q$ は可換であり、$\mathbb C(p,q)=\mathbb C(c,r)$、
逆変換は $p=r^2/c$、$q=r^3/c$ である。
導分 $2p\partial_p+3q\partial_q$ は $r\partial_r$ になるから、
$[P,r]=h$ かつ $c$ は中心的である。逆に $s=rP$ により全分数環が復元される。
標数零の Weyl 分数環の中心はその係数体に一致するので、最後の主張も従う。
\end{proof}

この構成の半古典極限は $\{P,r\}=1$、$c$ は Casimir である。
従って例外成分の一般元は、$\mathbb P^1_c\times\mathbb P^2$ の
三次式 $X_0^3$ に対応する垂直型の量子化と非可換双有理同値になる。
これは垂直成分内の特別な退化との対応である。

\subsection{Sklyanin pencil の中心部分体を拡大した比較}
\label{subsec:sklyanin-central-extension}

$S_4=S_4(E,\alpha)$ を四次元 Sklyanin 代数とし、
$\alpha$ は無限位数の平行移動とする。
以下の中央二次式と quadric pencil の性質は
\cite[\S10.1--10.4、定理10.2]{SmithVdB} による。
$\Omega_1,\Omega_2\in Z(S_4)_2$ を一次独立に取り、
\[
 u=\Omega_1\Omega_2^{-1},\qquad K=\mathbb C(u),\qquad
 Q_u=(S_4\otimes_{\mathbb C}K)/(\Omega_1-u\Omega_2)
\]
と置く。$u$ は超越的であり、$K$ は $Q_0(S_4)$ の中心部分体である。
実際、二つの中央二次式の同次な代数関係があれば、
複素数上の二変数斉次式の分解と domain 性から一次従属を生じる。
$Q_0(Q_u)\simeq Q_0(S_4)$ が成り立つ。
実際、$\Omega_2$ を逆にすると $u=\Omega_1\Omega_2^{-1}$ による同一視ができ、
残りは中心的な非零多項式の局所化になる。

\begin{proposition}\label{prop:sklyanin-finite-base-extension}
上の設定で、有限拡大 $K'/K$ と、$\mathbb C$ 上で定義された平行移動
$\sigma$ で $\sigma^3=\alpha^2$ を満たすものが存在して
\[
 Q_0(Q_u\otimes_KK')
 \simeq
 Q_0(S_3(E,\sigma)\otimes_{\mathbb C}K')
\]
が成り立つ。
\end{proposition}
\begin{proof}
$\overline K$ 上で $Q_u$ は滑らかな Van den Bergh quadric である。
これは pencil の特異な quadric が四つであることから従う
\cite[定理10.2]{SmithVdB}。
Rogalski--Sierra--Stafford \cite[定理1.2、定理7.3]{RSS2023} により、
$Q_u^{(2)}$ の一点ブローアップは、ある $S_3(E,\sigma)^{(3)}$ の二点ブローアップと同型であり、
$\sigma^3=\alpha^2$ である。従って $\overline K$ 上で表示した二つの関数体は同型になる。
$\alpha$ は定数で、$[3]\colon E\to E$ のファイバーは $\mathbb C$ 上有限であるから、
$\sigma$ は $\mathbb C$ 上で選べる。次数3の直線束の選択も、$\overline K$ 上の平行移動により
固定したものに置き換えられる。

両側の $2$ 次および $3$ 次 Veronese は次数1で生成され、次数1の中心元を持つ。
$Q_u$ 側ではその中心元を $z=\Omega_2$ と取り、
$(Q_u^{(2)}[z^{-1}])_0$ は $Q_u^{(2)}$ の次数1の生成元を $z$ で割った元で生成される。
$S_3^{(3)}$ 側も次数3の中心元で同様である。
従って両方の次数零部分は有限生成であり、それぞれの分数環が対象の関数体となる。
上の同型と逆同型による有限個の生成元の像を分数で書くと、現れる係数は有限個である。
それらは $K$ 上代数的なので、全てを含む有限拡大 $K'/K$ を取れる。
同型と逆同型の等式もその拡大上で成立し、結論を得る。
\end{proof}

命題\ref{prop:sklyanin-finite-base-extension}は中心部分体の有限拡大後の主張である。
$K$ 上の同型を得るための下降は残る課題である。
特に二つの ruling の交換と、得られる三次元 Sklyanin 代数の形式を同時に調べる必要がある。
固定した $\Omega$ の各ファイバーの比較を $K$ 上で同時化することが、ここでの下降問題である。
